% 04-monitor-control.tex - Section 4: Monitor and Control Interface

\chapter{Monitor and Control Interface}
\label{sec:monitor-control}

\section{Overview}

The NDP subsystem communicates with the Monitor and Control System (MCS) via UDP messages. Commands are received on port 1742 and responses are sent to port 1743.

\section{Message Format}

\subsection{Command Message Format}

Commands from MCS to NDP follow the standard LWA message format (Table~\ref{tab:cmd-format}).

\begin{table}[htbp]
    \centering
    \caption{Command Message Format}
    \label{tab:cmd-format}
    \begin{tabular}{llll}
        \toprule
        \textbf{Field} & \textbf{Bytes} & \textbf{Position} & \textbf{Description} \\
        \midrule
        Destination & 3 & 0--2 & Subsystem ID (``NDP'') \\
        Sender & 3 & 3--5 & Sender ID (``MCS'') \\
        Command & 3 & 6--8 & Command code \\
        Reference & 9 & 9--17 & Reference number \\
        Data Length & 4 & 18--21 & Length of data section \\
        MJD & 6 & 22--27 & Modified Julian Date \\
        MPM & 9 & 28--36 & Milliseconds Past Midnight \\
        - & 1 & 37 & Space \\
        Data & Variable & 38+ & Command-specific data \\
        \bottomrule
    \end{tabular}
\end{table}

\subsection{Example Command Message}

A \cmd{DRX} command from MCS to NDP:
\begin{verbatim}
NDP MCS DRX       123    15 60710  43200000 <binary data>
\end{verbatim}

Where:
\begin{itemize}
    \item Destination = ``NDP'' (padded to 3 chars)
    \item Sender = ``MCS'' (padded to 3 chars)
    \item Command = ``DRX'' (padded to 3 chars)
    \item Reference = 123 (right-justified in 9 chars)
    \item Data Length = 15 (right-justified in 4 chars)
    \item MJD = 60710 (right-justified in 6 chars)
    \item MPM = 43200000 (right-justified in 9 chars)
    \item Followed by 15 bytes of binary command data
\end{itemize}

\subsection{Response Message Format}

Responses from NDP to MCS include the command status and any requested data (Table~\ref{tab:resp-format}).

\begin{table}[htbp]
    \centering
    \caption{Response Message Format}
    \label{tab:resp-format}
    \begin{tabular}{llll}
        \toprule
        \textbf{Field} & \textbf{Bytes} & \textbf{Position} & \textbf{Description} \\
        \midrule
        Destination & 3 & 0--2 & Destination ID (echo of sender) \\
        Sender & 3 & 3--5 & ``NDP'' \\
        Command & 3 & 6--8 & Echo of command \\
        Reference & 9 & 9--17 & Echo of reference number \\
        Data Length & 4 & 18--21 & Length of data section \\
        MJD & 6 & 22--27 & Response MJD \\
        MPM & 9 & 28--36 & Response MPM \\
        - & 1 & 37 & Space \\
        Accept/Reject & 1 & 38 & ``A'' (accept) or ``R'' (reject) \\
        Status & 7 & 39--45 & Status code \\
        Data & Variable & 46+ & Response data \\
        \bottomrule
    \end{tabular}
\end{table}

\section{Management Information Base (MIB)}
\label{sec:mib}

The following MIB entries can be queried using the \cmd{RPT} command.

\subsection{General Information}

Table~\ref{tab:mib-general} lists the general MIB entries that identify the subsystem and report its overall status.

\begin{table}[htbp]
    \centering
    \caption{General MIB Entries}
    \label{tab:mib-general}
    \begin{tabular}{llp{6cm}}
        \toprule
        \textbf{MIB Entry} & \textbf{Type} & \textbf{Description} \\
        \midrule
        \mib{SUMMARY} & char[7] & Current system status \\
        \mib{INFO} & char[256] & Detailed status information \\
        \mib{LASTLOG} & char[256] & Last log entry \\
        \mib{SUBSYSTEM} & char[3] & Subsystem identifier (``NDP'') \\
        \mib{SERIALNO} & char[5] & Serial number \\
        \mib{VERSION} & char[256] & Software version \\
        \bottomrule
    \end{tabular}
\end{table}

\subsection{System Configuration}

Table~\ref{tab:mib-config} lists the fixed system configuration parameters.

\begin{table}[htbp]
    \centering
    \caption{System Configuration MIB Entries}
    \label{tab:mib-config}
    \begin{tabular}{llp{6cm}}
        \toprule
        \textbf{MIB Entry} & \textbf{Type} & \textbf{Description} \\
        \midrule
        \mib{NUM\_STANDS} & uint16 & Number of stands (256) \\
        \mib{NUM\_SERVERS} & uint8 & Number of servers (4) \\
        \mib{NUM\_BOARDS} & uint8 & Number of F-engine boards (16) \\
        \mib{NUM\_BEAMS} & uint8 & Number of beams (4) \\
        \mib{NUM\_DRX\_TUNINGS} & uint8 & Tunings per beam (2) \\
        \mib{NUM\_FREQ\_CHANS} & uint16 & Number of frequency channels (4096) \\
        \mib{STAT\_SAMPLE\_SIZE} & uint32 & Samples used for statistics (512) \\
        \bottomrule
    \end{tabular}
\end{table}

\subsection{DRX Configuration}

Table~\ref{tab:mib-drx} lists the per-tuning DRX configuration entries, where \texttt{\{n\}} ranges over active tunings.

\begin{table}[htbp]
    \centering
    \caption{DRX Configuration MIB Entries}
    \label{tab:mib-drx}
    \begin{tabular}{llp{6cm}}
        \toprule
        \textbf{MIB Entry} & \textbf{Type} & \textbf{Description} \\
        \midrule
        \mib{DRX\{n\}\_CONFIG\_FREQ} & float64 & Tuning \{n\} frequency in Hz \\
        \mib{DRX\{n\}\_CONFIG\_FILTER} & uint16 & Tuning \{n\} filter code \\
        \mib{DRX\{n\}\_CONFIG\_GAIN} & uint16 & Tuning \{n\} gain setting \\
        \mib{DRX\{n\}\_CONFIG\_HIGH\_DR} & uint8 & Tuning \{n\} high DR mode \\
        \bottomrule
    \end{tabular}
\end{table}

\subsection{TBS Configuration}

Table~\ref{tab:mib-tbs} lists the TBS configuration entries.

\begin{table}[htbp]
    \centering
    \caption{TBS Configuration MIB Entries}
    \label{tab:mib-tbs}
    \begin{tabular}{llp{6cm}}
        \toprule
        \textbf{MIB Entry} & \textbf{Type} & \textbf{Description} \\
        \midrule
        \mib{TBS\_CONFIG\_FREQ} & float64 & TBS center frequency in Hz \\
        \mib{TBS\_CONFIG\_FILTER} & uint16 & TBS filter code \\
        \bottomrule
    \end{tabular}
\end{table}

\subsection{Input Statistics}

Table~\ref{tab:mib-ant} lists the per-input statistics entries, where \texttt{\{n\}} is the 1-based input index (stand $\times$ polarization).

\begin{table}[htbp]
    \centering
    \caption{Input Statistics MIB Entries}
    \label{tab:mib-ant}
    \begin{tabular}{llp{6cm}}
        \toprule
        \textbf{MIB Entry} & \textbf{Type} & \textbf{Description} \\
        \midrule
        \mib{ANT\{n\}\_RMS} & float32 & RMS of samples for input \{n\} \\
        \mib{ANT\{n\}\_DCOFFSET} & float32 & Mean (DC offset) of samples \\
        \mib{ANT\{n\}\_SAT} & int32 & Saturated sample count \\
        \mib{ANT\{n\}\_PEAK} & int32 & Peak sample value \\
        \bottomrule
    \end{tabular}
\end{table}

\subsection{F-engine Board Status}

Table~\ref{tab:mib-board} lists the per-board F-engine status entries, where \texttt{\{n\}} is the 1-based board index.

\begin{table}[htbp]
    \centering
    \caption{F-engine Board MIB Entries}
    \label{tab:mib-board}
    \begin{tabular}{llp{6cm}}
        \toprule
        \textbf{MIB Entry} & \textbf{Type} & \textbf{Description} \\
        \midrule
        \mib{BOARD\{n\}\_STAT} & char[256] & Board \{n\} status (JSON) \\
        \mib{BOARD\{n\}\_INFO} & char[256] & Board \{n\} information (JSON) \\
        \mib{BOARD\{n\}\_TEMP\_MAX} & float32 & Maximum temperature \\
        \mib{BOARD\{n\}\_TEMP\_MIN} & float32 & Minimum temperature \\
        \mib{BOARD\{n\}\_TEMP\_AVG} & float32 & Average temperature \\
        \mib{BOARD\{n\}\_FIRMWARE} & char[256] & Firmware version \\
        \mib{BOARD\{n\}\_HOSTNAME} & char[256] & Network hostname \\
        \bottomrule
    \end{tabular}
\end{table}

\subsection{X-engine Server Status}

Table~\ref{tab:mib-server} lists the per-server X-engine status entries, where \texttt{\{n\}} is the 1-based server index.

\begin{table}[htbp]
    \centering
    \caption{X-engine Server MIB Entries}
    \label{tab:mib-server}
    \begin{tabular}{llp{6cm}}
        \toprule
        \textbf{MIB Entry} & \textbf{Type} & \textbf{Description} \\
        \midrule
        \mib{SERVER\{n\}\_STAT} & uint64 & Server \{n\} status \\
        \mib{SERVER\{n\}\_INFO} & char[256] & Server \{n\} information \\
        \mib{SERVER\{n\}\_SOFTWARE} & char[256] & Software version \\
        \mib{SERVER\{n\}\_TEMP\_MAX} & float32 & Maximum temperature \\
        \mib{SERVER\{n\}\_TEMP\_MIN} & float32 & Minimum temperature \\
        \mib{SERVER\{n\}\_TEMP\_AVG} & float32 & Average temperature \\
        \mib{SERVER\{n\}\_HOSTNAME} & char[256] & Network hostname \\
        \bottomrule
    \end{tabular}
\end{table}

\begin{calloutBox}{Implementation Note}{blue}
Server index 1 refers to the headnode. Indices 2--5 refer to the servers running the X-engine pipelines.
\end{calloutBox}

\subsection{Global Status}

Table~\ref{tab:mib-global} lists system-wide status entries.

\begin{table}[htbp]
    \centering
    \caption{Global Status MIB Entries}
    \label{tab:mib-global}
    \begin{tabular}{llp{6cm}}
        \toprule
        \textbf{MIB Entry} & \textbf{Type} & \textbf{Description} \\
        \midrule
        \mib{GLOBAL\_TEMP\_MAX} & float32 & Maximum system temperature \\
        \mib{GLOBAL\_TEMP\_MIN} & float32 & Minimum system temperature \\
        \mib{GLOBAL\_TEMP\_AVG} & float32 & Average system temperature \\
        \mib{CLK\_VAL} & uint32 & Milliseconds past midnight \\
        \mib{UTC\_START} & char[256] & UTC start time string \\
        \mib{UPTIME} & uint32 & System uptime in seconds \\
        \bottomrule
    \end{tabular}
\end{table}

\subsection{Command Status}

Table~\ref{tab:mib-cmdstat} lists the command status MIB entry.

\begin{table}[htbp]
    \centering
    \caption{Command Status MIB Entry}
    \label{tab:mib-cmdstat}
    \begin{tabular}{llp{6cm}}
        \toprule
        \textbf{MIB Entry} & \textbf{Type} & \textbf{Description} \\
        \midrule
        \mib{CMD\_STAT} & struct & Command status for previous slot \\
        \bottomrule
    \end{tabular}
\end{table}

The \mib{CMD\_STAT} structure contains:
\begin{itemize}
    \item \texttt{slot\_time} (uint32): Seconds past station time midnight
    \item \texttt{num\_commands} (uint16): Number of commands ($N \leq 606$)
    \item \texttt{reference[N]} (uint32): Reference numbers
    \item \texttt{completion\_code[N]} (uint8): Exit codes
\end{itemize}

\section{Control Commands}
\label{sec:commands}

\subsection{PNG --- Ping}

The \cmd{PNG} command tests connectivity with the NDP subsystem (Table~\ref{tab:cmd-png}).

\begin{table}[htbp]
    \centering
    \caption{PNG Command}
    \label{tab:cmd-png}
    \begin{tabular}{ll}
        \toprule
        \textbf{Field} & \textbf{Value} \\
        \midrule
        Command & \cmd{PNG} \\
        Data & (none) \\
        Response & Status only \\
        \bottomrule
    \end{tabular}
\end{table}

\subsection{INI --- Initialize}

The \cmd{INI} command initializes the NDP subsystem (Table~\ref{tab:cmd-ini}).

\begin{table}[htbp]
    \centering
    \caption{INI Command}
    \label{tab:cmd-ini}
    \begin{tabular}{ll}
        \toprule
        \textbf{Field} & \textbf{Value} \\
        \midrule
        Command & \cmd{INI} \\
        Data & Mode string (optional flags) \\
        Response & Status; on failure: exit code and message \\
        \bottomrule
    \end{tabular}
\end{table}

The initialization sequence:
\begin{enumerate}
    \item Verify F-engine boards are accessible
    \item Program FPGA firmware on all ZCU102 boards
    \item Configure F-engine parameters (FFT shift, equalization)
    \item Verify X-engine servers are accessible
    \item Start X-engine pipelines
    \item Verify data flow from F-engines to X-engines
\end{enumerate}

\subsubsection{Exit Codes}

\begin{tabular}{cl}
    \hex{00} & Process accepted without error \\
    \hex{04} & Board programming failed \\
    \hex{05} & Board configuration failed \\
    \hex{08} & Board-level shutdown failed \\
    \hex{09} & Server startup failed \\
    \hex{0A} & Server shutdown failed \\
    \hex{0B} & Pipeline startup failed \\
    \hex{0C} & Blocking operation in progress \\
    \hex{0D} & Internal F-engine error condition \\
    \hex{0E} & Pipeline processing error \\
\end{tabular}

\subsubsection{Options}

The \cmd{INI} command accepts the option strings listed in Table~\ref{tab:ini-options} (space-separated).

\begin{table}[htbp]
    \centering
    \caption{INI Command Options}
    \label{tab:ini-options}
    \begin{tabular}{lp{9cm}}
        \toprule
        \textbf{Option} & \textbf{Description} \\
        \midrule
        \texttt{FORCE} & Continue initialization even if errors occur on some boards \\
        \texttt{NOREPROGRAM} & Skip FPGA programming (assume already programmed) \\
        \texttt{STARTFREQ\_NNN} & Set the starting frequency of the observed band to \texttt{NNN}~MHz. The value is rounded to the nearest multiple of 16 channels ($\approx$383~kHz). If not specified, the default starting frequency from the configuration file is used. \\
        \bottomrule
    \end{tabular}
\end{table}

\subsection{SHT --- Shutdown}

The \cmd{SHT} command shuts down the NDP subsystem (Table~\ref{tab:cmd-sht}).

\begin{table}[htbp]
    \centering
    \caption{SHT Command}
    \label{tab:cmd-sht}
    \begin{tabular}{ll}
        \toprule
        \textbf{Field} & \textbf{Value} \\
        \midrule
        Command & \cmd{SHT} \\
        Data & Mode string (optional, space-separated) \\
        Response & Status; on failure: exit code and message \\
        \bottomrule
    \end{tabular}
\end{table}

\subsubsection{Exit Codes}

\begin{tabular}{cl}
    \hex{00} & Process accepted without error \\
    \hex{07} & Invalid command arguments (unknown mode) \\
    \hex{08} & Board-level shutdown failed \\
    \hex{0A} & Server shutdown failed \\
    \hex{0C} & Blocking operation in progress \\
\end{tabular}

\subsubsection{Options}

The available options are listed in Table~\ref{tab:sht-options}.

\begin{table}[htbp]
    \centering
    \caption{SHT Command Options}
    \label{tab:sht-options}
    \begin{tabular}{lp{9cm}}
        \toprule
        \textbf{Option} & \textbf{Description} \\
        \midrule
        \texttt{SCRAM} & Immediate hard shutdown without cleanup \\
        \texttt{RESTART} & Restart system after shutdown \\
        \texttt{FORCE} & Continue shutdown even if errors occur \\
        \texttt{HARD} & Reboot F-engine boards instead of just unprogramming \\
        \bottomrule
    \end{tabular}
\end{table}

\subsection{STP --- Stop}

The \cmd{STP} command stops an active observing mode (Table~\ref{tab:cmd-stp}).

\begin{table}[htbp]
    \centering
    \caption{STP Command}
    \label{tab:cmd-stp}
    \begin{tabular}{ll}
        \toprule
        \textbf{Field} & \textbf{Value} \\
        \midrule
        Command & \cmd{STP} \\
        Data & Mode to stop: ``DRX'', ``TBS'', ``TBT'', ``BEAM\{n\}'', ``COR'' \\
        Response & Status; on failure: exit code and message \\
        \bottomrule
    \end{tabular}
\end{table}

\subsubsection{Exit Codes}

\begin{tabular}{cl}
    \hex{00} & Process accepted without error \\
    \hex{63} & Subsystem not ready \\
    \hex{FF} & Invalid mode string \\
\end{tabular}

\subsubsection{Mode Strings}

The valid mode strings are listed in Table~\ref{tab:stp-modes}.

\begin{table}[htbp]
    \centering
    \caption{STP Mode Strings}
    \label{tab:stp-modes}
    \begin{tabular}{lp{9cm}}
        \toprule
        \textbf{Mode} & \textbf{Description} \\
        \midrule
        \texttt{DRX} & Stop all DRX/beam outputs \\
        \texttt{TBS} & Stop TBS streaming mode \\
        \texttt{TBT} & Stop TBT capture \\
        \texttt{BEAM1} -- \texttt{BEAM4} & Stop specific beam output \\
        \texttt{COR} & Stop correlator output \\
        \bottomrule
    \end{tabular}
\end{table}

\begin{calloutBox}{Implementation Note}{blue}
The \texttt{TBT} and \texttt{COR} stop modes are currently no-ops.  TBT captures complete on their own and COR runs continuously once the system is initialized.
\end{calloutBox}

\subsection{DRX --- Configure Beam Tuning}

The \cmd{DRX} command configures a frequency tuning for a beam (Table~\ref{tab:cmd-drx}).  The command data format is given in Table~\ref{tab:drx-data}.

\begin{table}[htbp]
    \centering
    \caption{DRX Command}
    \label{tab:cmd-drx}
    \begin{tabular}{ll}
        \toprule
        \textbf{Field} & \textbf{Value} \\
        \midrule
        Command & \cmd{DRX} \\
        Data & Binary packed structure (see below) \\
        Response & Status; on failure: exit code and message \\
        \bottomrule
    \end{tabular}
\end{table}

\begin{table}[htbp]
    \centering
    \caption{DRX Command Data Format (Big-Endian)}
    \label{tab:drx-data}
    \begin{tabular}{llll}
        \toprule
        \textbf{Field} & \textbf{Type} & \textbf{Range} & \textbf{Description} \\
        \midrule
        beam & uint8 & 1--4 & Beam number \\
        tuning & uint8 & 1--2 & Tuning number within beam \\
        freq & float64 & 0--98~MHz & Center frequency in Hz \\
        filter & uint8 & 1--7 & Filter code (see Table~\ref{tab:drx-filters}) \\
        gain & int16 & 0--15 & Gain setting (right bit-shift) \\
        high\_dr & uint8 & 0--1 & High dynamic range mode \\
        subslot & uint8 & 0--99 & Subslot for execution \\
        \bottomrule
    \end{tabular}
\end{table}

Python struct format: \texttt{'>BBdBhBB'} (15 bytes total)

\begin{table}[htbp]
    \centering
    \caption{DRX Filter Codes}
    \label{tab:drx-filters}
    \begin{tabular}{cl}
        \toprule
        \textbf{Code} & \textbf{Sample Rate} \\
        \midrule
        1 & 250~kHz \\
        2 & 500~kHz \\
        3 & 1.0~MHz \\
        4 & 2.0~MHz \\
        5 & 4.9~MHz \\
        6 & 9.8~MHz \\
        7 & 19.6~MHz \\
        \bottomrule
    \end{tabular}
\end{table}

\subsubsection{Exit Codes}

\begin{tabular}{cl}
    \hex{00} & Process accepted without error \\
    \hex{63} & Subsystem not ready \\
    \hex{FF} & Invalid parameter \\
\end{tabular}

\subsection{TBS --- Transient Buffer -- Streaming}

The \cmd{TBS} command configures the triggered buffer spectra mode (Table~\ref{tab:cmd-tbs}).  The command data format is given in Table~\ref{tab:tbs-data}.

\begin{table}[htbp]
    \centering
    \caption{TBS Command}
    \label{tab:cmd-tbs}
    \begin{tabular}{ll}
        \toprule
        \textbf{Field} & \textbf{Value} \\
        \midrule
        Command & \cmd{TBS} \\
        Data & Binary packed structure (see below) \\
        Response & Status; on failure: exit code and message \\
        \bottomrule
    \end{tabular}
\end{table}

\begin{table}[htbp]
    \centering
    \caption{TBS Command Data Format (Big-Endian)}
    \label{tab:tbs-data}
    \begin{tabular}{llll}
        \toprule
        \textbf{Field} & \textbf{Type} & \textbf{Range} & \textbf{Description} \\
        \midrule
        freq & float64 & 0--98~MHz & Center frequency in Hz \\
        filter & uint8 & 7--9 & Filter code (see Table~\ref{tab:tbs-filters})\\
        \bottomrule
    \end{tabular}
\end{table}

Python struct format: \texttt{'>dB'} (9 bytes total)

\begin{calloutBox}{Implementation Note}{blue}
TBS configuration commands are sent to all X-engine pipelines but a pipeline will only output data if it handles the entire frequency range requested.  Because of this any TBS requests that falls on a X-engine server/pipeline boundary will be accepted but not produce data.
\end{calloutBox}

\begin{table}[htbp]
    \centering
    \caption{TBS Filter Codes}
    \label{tab:tbs-filters}
    \begin{tabular}{cl}
        \toprule
        \textbf{Code} & \textbf{Sample Rate} \\
        \midrule
        7 & $\approx$95.7~kHz \\
        8 & $\approx$191.4~kHz \\
        9 & $\approx$287.1~kHz \\
        \bottomrule
    \end{tabular}
\end{table}

\subsection{TBT --- Transient Buffer -- Triggered}

The \cmd{TBT} command configures and triggers a time-domain voltage capture (Table~\ref{tab:cmd-tbt}).  The command data format is given in Table~\ref{tab:tbt-data}.

\begin{table}[htbp]
    \centering
    \caption{TBT Command}
    \label{tab:cmd-tbt}
    \begin{tabular}{ll}
        \toprule
        \textbf{Field} & \textbf{Value} \\
        \midrule
        Command & \cmd{TBT} \\
        Data & Binary packed structure (see below) \\
        Response & Status; on failure: exit code and message \\
        \bottomrule
    \end{tabular}
\end{table}

\begin{table}[htbp]
    \centering
    \caption{TBT Command Data Format (Big-Endian)}
    \label{tab:tbt-data}
    \begin{tabular}{llll}
        \toprule
        \textbf{Field} & \textbf{Type} & \textbf{Range} & \textbf{Description} \\
        \midrule
        trigger & uint64 & Full range & Trigger time in \fs{} ticks \\
        samples & int32 & Full range & Number of samples to capture \\
        mask & int64 & Full range & Stand mask (negative = local dump) \\
        \bottomrule
    \end{tabular}
\end{table}

Python struct format: \texttt{'>Qiq'} (20 bytes total)

\begin{calloutBox}{Implementation Note}{blue}
If the mask value is negative, the absolute value is used as the mask and the data is written to local storage rather than transmitted over the network.
\end{calloutBox}

\subsection{BAM --- Beamformer Coefficients}

The \cmd{BAM} command sets the beamformer delays and gains for a beam (Table~\ref{tab:cmd-bam}).  The command data format is given in Table~\ref{tab:bam-data}.

\begin{table}[htbp]
    \centering
    \caption{BAM Command}
    \label{tab:cmd-bam}
    \begin{tabular}{ll}
        \toprule
        \textbf{Field} & \textbf{Value} \\
        \midrule
        Command & \cmd{BAM} \\
        Data & Binary packed structure (see below) \\
        Response & Status; on failure: exit code and message \\
        \bottomrule
    \end{tabular}
\end{table}

\begin{table}[htbp]
    \centering
    \caption{BAM Command Data Format (Big-Endian)}
    \label{tab:bam-data}
    \begin{tabular}{llll}
        \toprule
        \textbf{Field} & \textbf{Type} & \textbf{Range} & \textbf{Description} \\
        \midrule
        beam & uint16 & 1--4 & Beam number \\
        delays & uint16[512] & Full range & Delay per input (fixed 16.4) \\
        gains & uint16[256][2][2] & Full range & Gain matrix per stand \\
        subslot & uint8 & 0--99 & Subslot for execution \\
        \bottomrule
    \end{tabular}
\end{table}

The delays array contains one 16-bit value per input (512 total). The value is in fixed-point format with 12 integer bits and 4 fractional bits, representing the delay in samples.

The gains array contains a $2 \times 2$ polarization mixing matrix for each stand. The matrix layout is:
\[
\begin{bmatrix}
G_{XX} & G_{XY} \\
G_{YX} & G_{YY}
\end{bmatrix}
\]

\subsubsection{Exit Codes}

\begin{tabular}{cl}
    \hex{00} & Process accepted without error \\
    \hex{63} & Subsystem not ready \\
    \hex{FF} & Invalid parameter \\
\end{tabular}

\subsection{COR --- Correlator Configuration}

The \cmd{COR} command configures the correlator (Table~\ref{tab:cmd-cor}).  The command data format is given in Table~\ref{tab:cor-data}.

\begin{table}[htbp]
    \centering
    \caption{COR Command}
    \label{tab:cmd-cor}
    \begin{tabular}{ll}
        \toprule
        \textbf{Field} & \textbf{Value} \\
        \midrule
        Command & \cmd{COR} \\
        Data & Binary packed structure (see below) \\
        Response & Status; on failure: exit code and message \\
        \bottomrule
    \end{tabular}
\end{table}

\begin{table}[htbp]
    \centering
    \caption{COR Command Data Format (Big-Endian)}
    \label{tab:cor-data}
    \begin{tabular}{llll}
        \toprule
        \textbf{Field} & \textbf{Type} & \textbf{Range} & \textbf{Description} \\
        \midrule
        navg & int32 & $\geq 500$ & Integration time in subslots \\
        gain & int16 & 0--15 & Gain setting (right bit-shift) \\
        subslot & uint8 & 0--99 & Subslot for execution \\
        \bottomrule
    \end{tabular}
\end{table}

Python struct format: \texttt{'>ihB'} (7 bytes total)

\subsubsection{Exit Codes}

\begin{tabular}{cl}
    \hex{00} & Process accepted without error \\
    \hex{63} & Subsystem not ready \\
    \hex{FF} & Invalid parameter \\
\end{tabular}

\begin{calloutBox}{Implementation Note}{blue}
The \texttt{COR} command is currently a no-op and the correlator output starts automatically with $\approx$5 s integrations once the system is initialized.
\end{calloutBox}

\subsection{RPT --- Report MIB Entry}

The \cmd{RPT} command queries a MIB entry (Table~\ref{tab:cmd-rpt}).

\begin{table}[htbp]
    \centering
    \caption{RPT Command}
    \label{tab:cmd-rpt}
    \begin{tabular}{ll}
        \toprule
        \textbf{Field} & \textbf{Value} \\
        \midrule
        Command & \cmd{RPT} \\
        Data & MIB entry name (ASCII string) \\
        Response & MIB value; on failure: error message \\
        \bottomrule
    \end{tabular}
\end{table}

\section{System States}

The NDP subsystem reports one of the status values listed in Table~\ref{tab:sys-status}.

\begin{table}[htbp]
    \centering
    \caption{System Status Values}
    \label{tab:sys-status}
    \begin{tabular}{lp{10cm}}
        \toprule
        \textbf{Status} & \textbf{Description} \\
        \midrule
        \texttt{SHUTDWN} & System is shut down or not yet initialized. No commands except \cmd{INI} will be processed. \\
        \texttt{BOOTING} & System is in the process of initializing. Commands will be rejected with exit code \hex{0C}. \\
        \texttt{NORMAL} & System is operating normally. All commands are accepted. \\
        \texttt{WARNING} & System is operating but a warning condition exists. Commands are accepted. \\
        \texttt{ERROR} & An error condition exists. The system may need to be re-initialized. \\
        \bottomrule
    \end{tabular}
\end{table}
